% Copyright (c) 2014,2016,2018 Casper Ti. Vector
% Public domain.

\chapter{方法}

\section{问题定义与总体训练范式}

\subsection{问题定义}

本文关注多模态大模型在开放世界场景下的输出安全性建模问题。具体而言，给定包含文本查询与视觉输入的多模态交互上下文，模型需要在保持有用性与指令遵循能力的同时，避免生成违反安全规范的内容，并具备对潜在风险的识别、分析与自我纠错能力。

形式化地，设多模态输入样本表示为：
\begin{equation}
x = (x_q, x_i)
\end{equation}
其中，$x_q$ 表示用户的文本查询，$x_i$ 表示与查询相关的图像输入（可为空）。模型的目标是生成文本回复 $y$，满足以下约束：

\begin{enumerate}
    \item \textbf{安全性（Safety）：} 输出 $y$ 不得违反预定义的安全规范，例如涉及违法指导、仇恨歧视、涉政敏感、暴力色情或其他高风险内容；
    \item \textbf{有用性（Helpfulness）：} 在满足安全约束前提下，输出应语义合理、信息充分，并尽可能满足用户真实意图。
\end{enumerate}

与传统仅在输出层面进行拒答或过滤的方法不同，本文强调模型在生成过程中具备显式的安全推理能力，即模型不仅能够判断“是否应当回答”，还应理解“为何该问题存在风险”以及“在风险约束下如何给出合理回应”。

\subsection{总体训练范式}

为实现上述目标，本文提出一种以规则约束与安全反思为核心的分阶段训练范式，通过将安全规范显式引入模型的推理与优化过程，使安全性由外部约束逐步内化为模型的生成行为。

整体训练框架包括三个关键组成部分：

\begin{enumerate}
    \item \textbf{基于规则的安全反思建模（Rule-grounded Safety Reflection）}

    针对不同类别的潜在安全风险，构建从高层安全原则（Principle）到实例级安全规则（Rule）的映射机制。规则以明确、可验证的形式描述模型在特定问题下应遵循的行为约束，并作为安全推理与训练监督的核心依据。

    \item \textbf{融合推理过程的有监督微调（Reasoning-aware SFT）}

    在监督学习阶段，通过合成包含安全规则、显式安全推理过程（Chain-of-Thought, CoT）以及最终回复的训练样本，引导模型在生成回复前进行安全反思与风险分析，建立“规则理解—风险推理—安全生成”的认知链条。

    \item \textbf{基于规则验证的在线强化学习（Rule-grounded Online RL, RLVR）}

    在初始安全能力建立后，引入在线强化学习流程。模型在 rollout 过程中生成候选回复，并通过基于规则验证的奖励机制对其安全性进行评估。奖励信号用于强化符合安全约束的生成行为，提升模型在复杂与分布外场景中的稳健性。
\end{enumerate}

本文采用如下两阶段训练策略：

\begin{itemize}
    \item \textbf{阶段一：安全反思监督微调（Safety Reflection SFT）}

    利用安全规则产生的思考与推理过程的合成数据，对模型进行有监督微调，使其具备基于规则进行安全分析与生成的能力的同时，隐式地将安全边界嵌入模型知识。

    \item \textbf{阶段二：规则验证驱动的在线强化学习（RLVR）}

    在阶段一模型基础上，引入在线强化学习，通过规则一致性验证提供奖励信号，进一步减少拒绝不足（Under-refusal）与越狱攻击成功率，同时避免显著损害模型的通用有用性。
\end{itemize}

该分阶段范式实现了安全推理能力从“被监督学习”到“行为层面内化”的逐步过渡。

\section{数据来源与预处理}
% 声明数据来源与数量, 预处理方式
在数据选择方面，本文选取了文小言线上产品及网页端所承接的真实用户多模态历史交互流量作为基础数据来源。该数据覆盖真实开放环境下的用户查询分布，包含文本与图像联合输入场景，同时也包含部分纯文本输入样本，能够较好反映实际部署环境中的安全风险形态与复杂程度。
为构建具有代表性的安全训练数据集，我们从原始流量中筛选出以下三类问题样本：

\begin{enumerate}
    \item \textbf{明确违规类别问题：} 已被初步判定为违规类别的问题，涵盖涉习涉政 血腥暴力 性感低俗 儿童邪典 等多种风险类型；
    \item \textbf{非违规常识问题：} 被判定为正常且属于常识性问答的问题，用于维持模型的基础有用性与正常响应能力, 缓解过度拒答现象；
    \item \textbf{边界型问题：} 分类结果为正常，但语义上接近违规边界、具有潜在风险或存在歧义的问题，用于增强模型在灰度场景下的判别与稳健能力。
\end{enumerate}

三类数据按照约 $1:10:2$ 的比例进行混合采样，最终构建规模约为 $4.6 \times 10^{4}$ 条的多模态训练数据集。其中包含文本-图像联合输入样本以及一定比例的纯文本样本，以保证模型在不同输入形态下的安全泛化能力。
\subsection{人工标注与质量控制}

为保证训练数据的准确性与一致性，我们对筛选后的数据进行了人工复核与精细化标注。具体而言：

\begin{itemize}
    \item 对样本是否属于违规内容进行人工确认；
    \item 对违规样本标注其具体违规类别；
    \item 对边界型问题进行风险等级与类别辅助标记；
    \item 对存在歧义或标注不一致的样本进行二次审核。
\end{itemize}

人工标注过程基于统一的安全分类体系与标注规范进行，确保不同类别之间定义清晰、边界明确。通过人工复核，有效降低了自动分类带来的误判与噪声，提高了后续规则构建与强化学习奖励建模的可靠性。

此外，所有数据在进入训练流程前均进行了脱敏处理，不包含可识别个人身份的信息，符合数据合规与隐私保护要求。

\section{基于规则的安全反思数据构建}

基于上述来源的多模态问题集，采用 Qwen3-VL-Instruct-8B 模型进行数据合成，以构建包含显式安全约束与推理过程的训练数据。

原始样本表示为三元组：
\begin{equation}
(x_q, x_c, x_i)
\end{equation}
其中，$x_q$ 为文本问题，$x_c$ 为类别标签，$x_i$ 为图像输入。

首先，根据类别 $x_c$ 构建类别到安全原则（Principle）的映射关系。随后结合具体问题 $x_q$，将抽象原则实例化为具体规则（Rule），用于限定模型生成行为。

数据合成采用两阶段流程：

\textbf{第一阶段：推理过程生成}
% 这里做了拒绝采样, 但是不知道怎么做的?

在给定 $x_q$、$x_i$ 及对应规则的条件下，引导模型按预设模板生成推理过程（Chain-of-Thought, CoT），并显式参考规则进行安全分析。

\textbf{第二阶段：最终回复生成}

在相同输入条件下，采用不同响应模板生成最终回答（Response），并保证其与第一阶段推理过程一致，同时满足规则约束。

最终每个样本扩展为三元结构：
\begin{equation}
(\text{Rule}, \text{CoT}, \text{Response})
\end{equation}

\subsection{基于规则验证的奖励建模}
% 参考了OpenAI的奖励建模设计, 需要指明!!
在强化学习阶段，本文设计了一个基于规则验证的多模态安全奖励验证器（Rule-grounded Reward Verifier）。该验证器不引入额外可学习参数，而是通过规则一致性判断与通用奖励模型相结合的方式，对模型生成结果进行分解式奖励评估。

具体而言，奖励信号由两个部分构成：安全性奖励 $r_{\text{safe}}$ 与有用性奖励 $r_{\text{help}}$，最终样本奖励定义为：

\begin{equation}
r = r_{\text{safe}} \cdot r_{\text{help}}
\end{equation}

\paragraph{(1) 安全违规类别样本}

对于被标注为明确违规类别的问题样本，奖励重点在于评估模型生成是否符合对应安全规则。

在 rollout 阶段，模型生成候选回复后，验证器依据该样本对应的规则（Rule）进行一致性校验，并进行二值打分：

\begin{equation}
r_{\text{safe}} =
\begin{cases}
1, & \text{若生成内容满足规则约束} \\
0, & \text{若生成内容违反规则}
\end{cases}
\end{equation}

对于此类样本，有用性奖励默认设为 $r_{\text{help}} = 1$，即仅通过安全一致性约束进行优化。该设计确保模型优先学习在高风险场景下的合规行为。

\paragraph{(2) 边界型与正常类别样本}

对于边界型问题及非违规常识问题，默认其安全性奖励为：

\begin{equation}
r_{\text{safe}} = 1
\end{equation}

此类样本的优化重点在于维持模型的指令遵循能力与生成质量。因此引入通用奖励模型（General Reward Model）对生成内容进行有用性评估，输出浮点数奖励：

\begin{equation}
r_{\text{help}} \in \mathbb{R}^{+}
\end{equation}

该奖励反映生成结果是否满足用户指令要求、语义完整性及回答质量。

\paragraph{(3) 奖励组合机制}

最终奖励值通过乘积形式融合安全性与有用性信号：

\begin{equation}
r = r_{\text{safe}} \cdot r_{\text{help}}
\end{equation}

该设计具有如下优势：

\begin{itemize}
    \item 当生成违反安全规则时，$r_{\text{safe}}=0$，整体奖励为零，从而对不安全行为进行强约束；
    \item 在满足安全约束前提下，通过 $r_{\text{help}}$ 进一步优化生成质量；
    \item 对不同类别样本采用差异化奖励策略，避免过度拒答并维持模型通用能力。
\end{itemize}

在 online-RL 训练过程中，模型对多模态输入进行 rollout 生成候选回复。验证器依据规则进行一致性与合规性校验，并结合通用奖励模型计算最终奖励值。
实现层面，基于自由RL框架, 将多模态输入统一转换为符合对话接口规范的格式，并采用异步 API 调用请求已部署好的教师通用模型获取奖励以提升并行吞吐效率，从而保证在线强化学习阶段的训练稳定性与计算效率。

奖励信号作为外部反馈参与参数更新，引导模型逐步学习符合安全规范的生成行为。


\section{Agentic RAG 场景下的安全验证探索}

除标准多模态输入场景外，本文进一步探索模型通过调用搜索工具, 解决一些事实性安全问题, 来提高安全性表现的方案。
具体而言，我们构建了一个面向 Agentic RAG 场景的事实性风险问题数据集，用于分析模型在“检索增强 + 安全判断”复合决策过程中的行为稳定性。

该数据集基于开源新闻标题构建。我们以新闻标题作为具有现实背景的查询起点，通过调用内部检索 API 召回相关文档，并选取相关性较高的文档作为补充证据。为保证文档可靠性, 仅选择了在信任域名单上的网页内容。
在此基础上，结合标题语境与召回内容自动生成具有潜在风险判断需求的事实性问题。这类问题通常涉及事实归因、事件评价或延伸推理，若缺乏时效性知识，模型容易产生不准确或具有误导性的输出。

与常规安全评测不同，该场景不仅要求模型识别内容风险，还需要其在不确定或信息不足时主动利用检索工具获取外部证据，并基于检索结果进行谨慎生成。因此，该实验主要关注以下能力：模型是否能够合理触发工具调用、是否基于证据进行推理，以及是否在事实不充分或存在争议时避免生成不当内容。
需要强调的是，本工作在 Agentic RAG 安全方向上的探索仍属于初步尝试。我们并未引入专门的工具规划或复杂决策优化机制，而是在现有安全训练框架下观察模型在外部检索条件下的行为变化。该实验旨在提供一个安全性与工具调用能力交叉评估的初步范式，为后续更系统的安全 Agent 研究提供经验参考。
